\documentclass[10pt,a4paper]{amsart}
\usepackage{amsmath,amssymb,amsthm,graphicx}
\usepackage{verbatim}
\newenvironment{reference}{\comment}{\endcomment}
%\newenvironment{reference}{}{}
\newenvironment{slogan}{\comment}{\endcomment}
\newenvironment{history}{\comment}{\endcomment}
% For commutative diagrams we use Xy-pic
\usepackage[all]{xy}

% We use 2cell for 2-commutative diagrams.
\xyoption{2cell}
% Theorem environments.
%
\theoremstyle{plain}
\newtheorem{theorem}[subsection]{Theorem}
\newtheorem{proposition}[subsection]{Proposition}
\newtheorem{lemma}[subsection]{Lemma}

\theoremstyle{definition}
\newtheorem{definition}[subsection]{Definition}
\newtheorem{example}[subsection]{Example}
\newtheorem{exercise}[subsection]{Exercise}
\newtheorem{situation}[subsection]{Situation}

\theoremstyle{remark}
\newtheorem{remark}[subsection]{Remark}
\newtheorem{remarks}[subsection]{Remarks}

\numberwithin{equation}{subsection}

% Macros
%
\def\lim{\mathop{\mathrm{lim}}\nolimits}
\def\colim{\mathop{\mathrm{colim}}\nolimits}
\def\Spec{\mathop{\mathrm{Spec}}}
\def\Hom{\mathop{\mathrm{Hom}}\nolimits}
\def\Ext{\mathop{\mathrm{Ext}}\nolimits}
\def\SheafHom{\mathop{\mathcal{H}\!\mathit{om}}\nolimits}
\def\SheafExt{\mathop{\mathcal{E}\!\mathit{xt}}\nolimits}
\def\Sch{\mathit{Sch}}
\def\Mor{\mathop{\mathrm{Mor}}\nolimits}
\def\Ob{\mathop{\mathrm{Ob}}\nolimits}
\def\Sh{\mathop{\mathit{Sh}}\nolimits}
\def\NL{\mathop{N\!L}\nolimits}
\def\CH{\mathop{\mathrm{CH}}\nolimits}
\def\proetale{{pro\text{-}\acute{e}tale}}
\def\etale{{\acute{e}tale}}
\def\QCoh{\mathit{QCoh}}
\def\Ker{\mathop{\mathrm{Ker}}}
\def\Im{\mathop{\mathrm{Im}}}
\def\Coker{\mathop{\mathrm{Coker}}}
\def\Coim{\mathop{\mathrm{Coim}}}

% Boxtimes
%
\DeclareMathSymbol{\boxtimes}{\mathbin}{AMSa}{"02}

%
% Macros for moduli stacks/spaces
%
\def\QCohstack{\mathcal{QC}\!\mathit{oh}}
\def\Cohstack{\mathcal{C}\!\mathit{oh}}
\def\Spacesstack{\mathcal{S}\!\mathit{paces}}
\def\Quotfunctor{\mathrm{Quot}}
\def\Hilbfunctor{\mathrm{Hilb}}
\def\Curvesstack{\mathcal{C}\!\mathit{urves}}
\def\Polarizedstack{\mathcal{P}\!\mathit{olarized}}
\def\Complexesstack{\mathcal{C}\!\mathit{omplexes}}
% \Pic is the operator that assigns to X its picard group, usage \Pic(X)
% \Picardstack_{X/B} denotes the Picard stack of X over B
% \Picardfunctor_{X/B} denotes the Picard functor of X over B
\def\Pic{\mathop{\mathrm{Pic}}\nolimits}
\def\Picardstack{\mathcal{P}\!\mathit{ic}}
\def\Picardfunctor{\mathrm{Pic}}
\def\Deformationcategory{\mathcal{D}\!\mathit{ef}}
\usepackage{cleveref}
\crefname{originallemma}{original lemma}{original lemmas}
\Crefname{originallemma}{Original lemma}{Original lemmas}
\makeatletter
\newlabel{lemma-constructible-hausdorff-quasi-compact}{{lemma-constructible-hausdorff-quasi-compact}{source}{}{sourcepointer-1}{}}
\newlabel{lemma-constructible-hausdorff-quasi-compact@cref}{{[originallemma][0][]lemma-constructible-hausdorff-quasi-compact}{[1][0][]source}{}{}{}}
\makeatother
\crefname{originallemma}{original lemma}{original lemmas}
\Crefname{originallemma}{Original lemma}{Original lemmas}
\makeatletter
\newlabel{lemma-fibres-spectral-map-quasi-compact}{{lemma-fibres-spectral-map-quasi-compact}{source}{}{sourcepointer-2}{}}
\newlabel{lemma-fibres-spectral-map-quasi-compact@cref}{{[originallemma][0][]lemma-fibres-spectral-map-quasi-compact}{[1][0][]source}{}{}{}}
\makeatother
\crefname{originallemma}{original lemma}{original lemmas}
\Crefname{originallemma}{Original lemma}{Original lemmas}
\makeatletter
\newlabel{lemma-image-quasi-compact}{{lemma-image-quasi-compact}{source}{}{sourcepointer-3}{}}
\newlabel{lemma-image-quasi-compact@cref}{{[originallemma][0][]lemma-image-quasi-compact}{[1][0][]source}{}{}{}}
\makeatother
\title{A map of spectral spaces is spectral exactly when it is continuous for the constructible topology}
\author{The Stacks Project Authors}
\date{stacks-project:topology.tex:local-snapshot-20261012; GNU FDL 1.2 or later}
\begin{document}
\maketitle
\noindent\textit{Editorial note.} All delivered mathematical text is the author's own, copied from the pinned Stacks Project chapter source sources/stacks-topology/pinned-original-source.tex; no statement or proof step was rewritten, reordered or removed.\par\medskip
\noindent\textit{Editorial note.} Preamble provenance: the Stacks Project keeps one single shared preamble for every chapter; this card takes it directly from the staged sources/stacks-topology/preamble.tex as three verbatim spans (lines 8-12, 14-18 and 154-217), so no derived file was written and nothing was added or redefined. The wrapper class and the base packages are supplied by the assembler.\par\medskip
\noindent\textit{Editorial note.} Bibliography provenance: neither the counted statement nor its proof contains a citation (this lane rescans the whole span before writing the card), so the card carries no bibliography and no entry was reconstructed.\par\medskip
\noindent\textit{Editorial note.} Selection provenance: this card was opened under the file-based control contract (CONTROL.md), with the mechanical pre-card gate run on its statement and proof interval before the card was written.\par\medskip
\noindent\textit{Allowed standard prerequisites (editorial).}\begin{itemize}
\item The statement that a spectral map is continuous in the constructible topology and has quasi-compact fibres. Invoked at the start of the counted proof (source line 4628). Gives the easy direction.
\item The statement that the constructible topology of a spectral space is Hausdorff and quasi-compact. Invoked inside the counted proof (source lines 4633-4635). Makes the preimage of a quasi-compact open quasi-compact in the constructible topology.
\item The statement that quasi-compactness transfers across a surjective continuous map. Invoked at the end of the counted proof (source line 4640). Transfers quasi-compactness back to the given topology.
\end{itemize}\medskip
\ifdefined\hypertarget\hypertarget{sourcepointer-1}{}\fi
\noindent\textit{Source reference.} Original lemma lemma-constructible-hausdorff-quasi-compact: sources/stacks-topology/pinned-original-source.tex, source lines 4522--4522. Navigation pointer only: the counted text refers to the Stacks item lemma-constructible-hausdorff-quasi-compact, which is an existing labeled result of this same chapter and is not part of this card. The pointer carries the target's own label as its number and its definition line as its span; its content is not reproduced here.\par\medskip
\ifdefined\hypertarget\hypertarget{sourcepointer-2}{}\fi
\noindent\textit{Source reference.} Original lemma lemma-fibres-spectral-map-quasi-compact: sources/stacks-topology/pinned-original-source.tex, source lines 4589--4589. Navigation pointer only: the counted text refers to the Stacks item lemma-fibres-spectral-map-quasi-compact, which is an existing labeled result of this same chapter and is not part of this card. The pointer carries the target's own label as its number and its definition line as its span; its content is not reproduced here.\par\medskip
\ifdefined\hypertarget\hypertarget{sourcepointer-3}{}\fi
\noindent\textit{Source reference.} Original lemma lemma-image-quasi-compact: sources/stacks-topology/pinned-original-source.tex, source lines 1964--1964. Navigation pointer only: the counted text refers to the Stacks item lemma-image-quasi-compact, which is an existing labeled result of this same chapter and is not part of this card. The pointer carries the target's own label as its number and its definition line as its span; its content is not reproduced here.\par\medskip
\begin{lemma}
\label{lemma-spectral-if-continuous-wrt-constructible-top}
Let $X$ and $Y$ be spectral spaces. Let $f : X \to Y$ be a continuous map.
Then $f$ is spectral if and only if $f$ is continuous in the constructible
topology.
\end{lemma}

\begin{proof}
The only if part of this is
Lemma \ref{lemma-fibres-spectral-map-quasi-compact}.
Assume $f$ is continuous in the constructible topology.
Let $V \subset Y$ be quasi-compact open.
Then $V$ is open and closed in the constructible topology.
Hence $f^{-1}(V)$ is open and closed in the constructible topology.
Hence $f^{-1}(V)$ is quasi-compact in the constructible topology
as $X$ is quasi-compact in the constructible topology by
Lemma \ref{lemma-constructible-hausdorff-quasi-compact}.
Since the identity $f^{-1}(V) \to f^{-1}(V)$ is surjective
and continuous from the constructible topology to the usual
topology, we conclude that $f^{-1}(V)$ is quasi-compact
in the topology of $X$ by Lemma \ref{lemma-image-quasi-compact}.
This finishes the proof.
\end{proof}

\end{document}
